\textbf{8.\ \textsc{Connected charges} (8 points)} --- \textit{Jaan Kalda.}

In the region $0 < x < L$, there is an electric field $\vec{E} = E_0\hat{x}$,
where $\hat{x}$ denotes the unit vector parallel to the $x$-axis. Two small
balls, each of mass $m$ and carrying charge $q$ ($qE_0 > 0$) are connected with
a weightless non-stretchable string of length $l$. Initially, at the moment of
time $t = 0$, the string is taut, the velocity of the both balls is
$\vec{v} = v_0\hat{x}$, one of the balls, the ball $A$, is at $x = 0$ while the
other ball, the ball $B$ is at $x < 0$. The electric field created by the balls
can be neglected, and it can be assumed that $v_0$ is very small (much smaller
than $\sqrt{ELq/m}$).

\textbf{i)} \textit{(2 points)}Consider the case when the string is parallel to
the $x$-axis, and $l = L$ Sketch the dependence of the velocity of the both
balls as a function of time. Will the balls collide? If yes then when?

\textbf{ii)} \textit{(2 points)} Now, at $t = 0$, the string forms an angle of
$45^\circ$ with the $x$-axis, $l = 1.2291L$. By this string length, at the
moment $t = T$ when the ball $A$ reaches $x = L$, the string is parallel to the
$x$-axis. Find $T$.

\textbf{iii)} \textit{(2 points)} Under the assumptions of the previous task,
what is the speed of the ball $A$ at the moment $t = T$?

\textbf{iv)} \textit{(2 points)} Now, at $t = 0$, the string forms a very small
angle $\phi$ with the $x$-axis. Similarly to the previous two tasks, at the
moment $t = T'$ when the ball $A$ reaches $x = L$, the string is parallel to the
$x$-axis. Find the string length $l$ assuming that $l > L$. \textit{The answer
to this point should contain only $L$ and numerical constants.}
